% !TeX root = ../main.tex

\xchapter{基于多模态融合的人体活动识别}{Crossreferences}

本章主要围绕基于多模态融合的人体活动识别任务展开深入研究。在掌握了相关理论与技术的基础上，第3.1节首先对当前人体活动识别领域面临的关键科学问题进行了系统性描述，重点分析了单一传感模态在复杂场景下的局限性，以及开展活动、身份与场景多任务协同识别的必要性。第3.2节详细介绍了实验所用的数据集构建过程，涵盖了数据采集的硬件配置、场景设置及数据预处理流程，为后续研究提供了坚实的数据支撑。第3.3节是本章的核心，全面阐述了本文提出的基于IMU和Wi-Fi多模态融合的识别方法。该部分详细解读了由IMU特征提取器、Wi-Fi特征提取器、交叉注意力融合模块以及多任务分类器组成的网络架构，并深入剖析了如何利用对比学习策略实现异构特征在同一语义空间的有效对齐，从而提升多任务识别的准确率。第3.4节则对实验设计与结果进行了详尽的分析，从评价指标、实验环境出发，通过数据实验、网络架构对比、消融实验以及跨域实验等多维度的评估，充分验证了所提方法在多任务识别中的有效性与鲁棒性。

\xsection{问题描述}{Ref labels}

人体活动识别（Human Activity Recognition, HAR）作为实现普适计算、环境智能以及人机自然交互的核心感知技术，致力于通过分析传感器数据来自动解析人类的行为模式与意图。近年来，随着微机电系统（MEMS）技术的成熟与无线通信基础设施的完善，基于智能穿戴设备中惯性测量单元（IMU）与环境泛在Wi-Fi信号的感知方案，凭借其硬件成本低廉、部署便捷、全天候可用以及具备非视距感知能力等优势，逐渐成为区别于计算机视觉感知的另一主流技术路线。特别是在家庭、医院等对隐私保护要求极高的私密空间中，这种非视觉感知方案展现出了巨大的应用潜力。然而，在追求高精度、高鲁棒性以及高语义密度的复杂场景感知目标时，现有的技术方案仍面临着多重深层次的科学挑战与技术瓶颈。

第一，单一传感模态在面对复杂动态环境时，受限于其物理感知机理，存在难以克服的鲁棒性瓶颈与感知盲区。
对于惯性测量单元（IMU）而言，虽然其内置的加速度计和陀螺仪能够以极高的采样率捕捉人体局部肢体的微观动力学特征（如震颤、挥动），但其本质上是一种本体感知传感器，缺乏对外部环境的感知能力。这意味着IMU难以区分那些动作模式相似但发生地点或环境交互不同的活动（例如，“在书房的打字行为”与“在厨房的切菜行为”在手腕IMU信号上可能表现出极高的相似性）。此外，受限于低成本MEMS传感器的制造工艺，IMU数据不可避免地包含高频噪声与零偏误差，在长时序的运动过程中，这些误差会随着积分运算累积，导致对运动轨迹的估计产生严重的漂移现象。
反观Wi-Fi感知技术，虽然其基于信道状态信息（CSI）能够捕捉人体在空间中运动所引起的电磁波反射、散射与衍射变化，从而实现非接触式的全身运动感知，但Wi-Fi信号对环境变化极度敏感。室内环境中的多径效应（Multipath Effect）使得CSI信号不仅包含人体运动特征，还混杂了大量来自墙壁、家具等静态物体的反射分量。当环境布局发生微小改变或存在非目标人员干扰时，信号的特征分布会发生显著偏移，导致模型的跨场景泛化能力大幅下降。此外，受限于Wi-Fi信号的波长与带宽，其空间分辨率相对有限，难以精准捕捉肢体末端的细微动作。
因此，如何突破单一模态的物理局限，通过多模态融合实现局部动力学特征与全局空间特征的优势互补，是提升系统鲁棒性的首要问题。

第二，异构传感器数据在特征空间中的语义对齐面临巨大的“模态鸿沟”（Modality Gap），导致传统融合方法难以发挥协同效应。
IMU数据与Wi-Fi CSI数据在物理属性、数据结构及采样频率上存在本质差异。IMU数据通常表现为本体坐标系下的多维时间序列，反映的是物体受力与旋转的运动学属性；而Wi-Fi CSI数据则表现为世界坐标系下频域上的复数矩阵，反映的是无线信道的多径传输特性。这种巨大的异构性使得两种数据在特征分布上存在显著的差异。
传统的融合方法大多采用数据层拼接或特征层简单的加权求和（Late Fusion），这种处理方式往往忽略了模态间的分布差异，难以挖掘出深层的互补关系。在缺乏有效对齐机制的情况下，神经网络往往会倾向于过拟合某一种信噪比更高的模态，而忽略另一种模态的有效信息，甚至可能引入噪声干扰，导致“1+1<2”的负迁移现象。因此，亟需一种先进的特征表示学习机制，能够跨越物理属性的差异，将IMU与Wi-Fi特征强制映射并对齐到同一个高维语义空间中，使得相同语义的动作在特征空间内距离更近，从而最大化多模态数据的协同感知能力。本文引入的对比学习（Contrastive Learning）策略正是为了解决这一关键的特征对齐难题。

第三，现有的人体活动识别模型大多局限于单一的“活动分类”任务，缺乏对身份、场景等关联语义的综合解析能力，无法满足复杂应用场景下的高层次智能化需求。
传统的HAR研究通常将问题简化为“识别当前发生了什么动作”，而忽略了在实际的物联网应用中，完整的语义上下文通常包含“谁（Identity）”、“在何处（Scene）”以及“做什么（Activity）”这三个维度的信息。例如，在智慧养老的跌倒检测场景中，系统不仅需要识别出“跌倒”这一动作，更需要确认跌倒者的身份（是老人还是护工）以及发生的地点（是高风险的浴室还是相对安全的卧室），这些综合信息对于触发精准的急救响应机制至关重要。
然而，现有的单任务模型通常将这些属性割裂开来独立训练，忽略了任务之间的内在耦合性。事实上，身份特征（如体型、步态）往往隐含在运动数据中，而场景特征（如Wi-Fi信号的多径指纹）也与活动识别共享底层的信号模式。若分别部署多个独立模型，不仅增加了系统的计算开销与延迟，也浪费了任务间共享特征带来的正则化增益。因此，本文旨在突破单一任务的限制，构建一个多任务协同学习（Multi-Task Learning）框架，通过共享底层的特征提取网络，并在高层设计并行的任务分类器，实现对“活动-身份-场景”全要素的一体化精准感知，从而显著提升系统的信息密度与应用价值。

综上所述，针对单一模态鲁棒性差、异构特征难对齐以及单任务信息量不足这三大核心挑战，本文提出了基于IMU与Wi-Fi多模态融合的人体活动识别方法。该方法通过对比学习跨越模态鸿沟，并通过多任务协同机制挖掘深层语义关联，旨在构建一个高精度、全方位的智能感知系统。

\xsection{数据集构建}{References}

\xsection{方法描述}{References}

\begin{figure}[h]
  \centering
  \includegraphics[width=\textwidth]{活动识别模型架构图.jpg}
  \caption{人体活动识别模型架构图}
\end{figure}

\xsection{实验与结果分析}{References}

\xsubsection{评价指标}{Ref equations}

\xsubsection{实验环境}{Level of this document}

\begin{figure}[h]
  \centering
  \includegraphics[width=\textwidth]{confusion_matrix.png}
  \caption{transformer}
\end{figure}

\xsubsection{数据实验结果分析}{Ref equations}

\subsubsection{模态消融实验}{Ref equations}

\begin{table}[H]
  \centering
  \begin{tblr}{
    colspec={X[c] X[c] X[c] X[c]},
  }
  \toprule
  模态 & 活动准确率 & 身份准确率 & 场景准确率 \\
  \midrule

  imu            & 92.96\% & 96.79\% & 97.83\% \\
  wifi           & 87.96\% & 99.70\% & 99.92\% \\
  wifi\_amp      & 87.75\% & 99.52\% & 99.89\% \\
  wifi\_pha      & 84.85\% & 97.86\% & 100.00\% \\
  stft           & 64.34\% & 96.63\% & 100.00\% \\
  stft\_amp      & 68.45\% & 97.54\% & 99.99\% \\
  stft\_phase    & 23.94\% & 63.80\% & 97.56\% \\
  imu + wifi     & 96.06\% & 99.93\% & 99.99\% \\
  imu + wifi\_amp & 95.68\% & 99.93\% & 99.95\% \\
  imu + wifi\_pha & 94.65\% & 99.70\% & 99.99\% \\
  imu + stft     & 94.00\% & 99.66\% & 99.92\% \\
  imu + stft\_amp & 93.39\% & 99.63\% & 99.88\% \\
  imu + stft\_pha & 90.24\% & 96.58\% & 98.78\% \\

  \bottomrule
  \end{tblr}
\end{table}

\subsubsection{设备消融实验}{Ref equations}

\begin{table}[H]
  \centering
  \begin{tblr}{
    colspec={X[c] X[c] X[c] X[c] X[c]},
  }
  \toprule
  设备 & 设备编号 & 活动准确率 & 身份准确率 & 场景准确率 \\
  \midrule

  -    & -    & 96.06\% & 99.93\% & 99.99\% \\
  imu  & 0    & 94.86\% & 99.85\% & 99.90\% \\
  imu  & 1    & 95.67\% & 99.94\% & 99.96\% \\
  imu  & 2    & 94.79\% & 99.82\% & 99.96\% \\
  imu  & 3    & 95.34\% & 99.86\% & 99.95\% \\
  imu  & 4    & 95.65\% & 99.88\% & 99.96\% \\
  wifi & 0    & 95.53\% & 99.76\% & 99.95\% \\
  wifi & 1    & 94.80\% & 99.77\% & 99.95\% \\
  wifi & 2    & 94.79\% & 99.71\% & 99.94\% \\

  \bottomrule
  \end{tblr}
\end{table}

\subsubsection{数据分割实验}{Ref equations}

\begin{table}[H]
  \centering
  \begin{tblr}{
    colspec={X[2,c] X[1,c] X[1,c] X[1,c]},
  }
  \toprule
  数据分割方案 & 活动准确率 & 身份准确率 & 场景准确率 \\
  \midrule

  window\_size = \ 25, stride = \ 25  & 88.83\% & 99.73\% & 99.95\% \\
  window\_size = \ 50, stride = \ 25  & 96.06\% & 99.93\% & 99.99\% \\
  window\_size = \ 75, stride = \ 25  & 98.50\% & 99.96\% & 99.99\% \\
  window\_size = 100, stride = \ 25  & 99.27\% & 99.96\% & 99.99\% \\
  window\_size = 125, stride = \ 25  & 99.62\% & 99.99\% & 99.99\% \\
  window\_size = 150, stride = \ 25  & 99.79\% & 99.98\% & 99.96\% \\
  window\_size = \ 50, stride = \ 50  & 88.71\% & 99.52\% & 99.90\% \\
  window\_size = \ 75, stride = \ 50  & 94.47\% & 99.85\% & 99.94\% \\
  window\_size = 100, stride = \ 50  & 96.95\% & 99.94\% & 99.93\% \\
  window\_size = 125, stride = \ 50  & 97.50\% & 99.68\% & 99.73\% \\
  window\_size = 150, stride = \ 50  & 98.77\% & 99.93\% & 99.90\% \\
  window\_size = \ 75, stride = \ 75  & 89.85\% & 99.54\% & 99.71\% \\
  window\_size = 100, stride = \ 75  & 91.82\% & 99.21\% & 99.85\% \\
  window\_size = 125, stride = \ 75  & 92.73\% & 98.38\% & 99.44\% \\
  window\_size = 150, stride = \ 75  & 96.33\% & 99.39\% & 99.72\% \\
  window\_size = 100, stride = 100 & 91.31\% & 99.12\% & 99.71\% \\
  window\_size = 125, stride = 100 & 92.34\% & 98.48\% & 99.30\% \\
  window\_size = 150, stride = 100 & 93.69\% & 98.40\% & 99.27\% \\
  window\_size = 125, stride = 125 & 90.87\% & 96.81\% & 99.00\% \\
  window\_size = 150, stride = 125 & 94.91\% & 99.28\% & 99.70\% \\
  window\_size = 150, stride = 150 & 93.21\% & 98.52\% & 99.30\% \\

  \bottomrule
  \end{tblr}
\end{table}

\xsubsection{网络架构实验结果分析}{Ref equations}

\subsubsection{骨干网络选择实验}{Ref equations}

\begin{table}[H]
  \centering
  \begin{tblr}{
    colspec={X[c] X[c] X[c] X[c]},
  }
  \toprule
  骨干网络 & 活动准确率 & 身份准确率 & 场景准确率 \\
  \midrule

  resnet18    & 88.44\% & 98.75\% & 99.99\% \\
  mamba       & 96.06\% & 99.93\% & 99.99\% \\
  lstm        & 91.43\% & 99.68\% & 99.91\% \\
  gru         & 92.90\% & 99.90\% & 99.99\% \\
  transformer & 95.06\% & 99.95\% & 99.97\% \\

  \bottomrule
  \end{tblr}
\end{table}

\subsubsection{模态融合方案选择实验}{Ref equations}

\begin{table}[H]
  \centering
  \begin{tblr}{
    colspec={X[c] X[c] X[c] X[c]},
  }
  \toprule
  模态融合方案 & 活动准确率 & 身份准确率 & 场景准确率 \\
  \midrule

  拼接       & 96.20\% & 99.84\% & 99.99\% \\
  交叉注意力  & 96.06\% & 99.93\% & 99.99\% \\

  \bottomrule
  \end{tblr}
\end{table}

\xsubsection{消融实验结果分析}{Ref equations}

\subsubsection{标准化消融实验}{Ref equations}

\begin{table}[H]
  \centering
  \begin{tblr}{
    colspec={X[c] X[c] X[c] X[c]},
  }
  \toprule
  IMU标准化 & WiFi标准化 & 活动准确率 & 身份准确率 & 场景准确率 \\
  \midrule

  否 & 否 & 96.06\% & 99.93\% & 99.99\% \\
  是 & 否 & 96.36\% & 99.94\% & 99.99\% \\
  否 & 是 & 95.60\% & 99.91\% & 99.99\% \\
  是 & 是 & 96.30\% & 99.94\% & 99.99\% \\

  \bottomrule
  \end{tblr}
\end{table}

\subsubsection{对比学习消融实验}{Ref equations}

\begin{table}[H]
  \centering
  \begin{tblr}{
    colspec={X[c] X[c] X[c] X[c]},
  }
  \toprule
  对比学习 & 活动准确率 & 身份准确率 & 场景准确率 \\
  \midrule

  否 & 96.06\% & 99.93\% & 99.99\% \\
  是 & 94.06\% & 99.75\% & 99.98\% \\

  \bottomrule
  \end{tblr}
\end{table}

\subsubsection{损失函数权重消融实验}{Ref equations}

\begin{table}[H]
  \centering
  \begin{tblr}{
    colspec={X[c] X[c] X[c] X[c]},
  }
  \toprule
  交叉熵损失函数权重 & 活动准确率 & 身份准确率 & 场景准确率 \\
  \midrule

  是 & 96.06\% & 99.93\% & 99.99\% \\
  否 & 94.76\% & 99.81\% & 99.92\% \\

  \bottomrule
  \end{tblr}
\end{table}

\xsubsection{跨域实验结果分析}{Ref equations}

\subsubsection{跨场景实验}{Ref equations}

\begin{table}[H]
  \centering
  \begin{tblr}{
    colspec={X[c] X[c] X[c] X[c]},
  }
  \toprule
  测试场景 & 活动准确率 & 身份准确率 & 场景准确率 \\
  \midrule

  livingroom & 25.58\% & 42.76\% & 3.94\% \\
  office     & 30.74\% & 34.71\% & 6.16\% \\
  kitchen    & 23.61\% & 4.62\% & 4.25\% \\

  \bottomrule
  \end{tblr}
\end{table}

\subsubsection{跨人实验}{Ref equations}

\begin{table}[H]
  \centering
  \begin{tblr}{
    colspec={X[c] X[c] X[c] X[c]},
  }
  \toprule
  测试志愿者 & 活动准确率 & 身份准确率 & 场景准确率 \\
  \midrule

  0  & 85.32\% & 99.39\% & 16.73\% \\
  1  & 55.14\% & 81.64\% & 11.48\% \\
  2  & 82.56\% & 99.29\% & 15.76\% \\
  3  & 74.60\% & 99.62\% & 1.20\% \\
  4  & 62.13\% & 98.56\% & 29.95\% \\
  5  & 76.29\% & 99.22\% & 0.13\% \\
  6  & 67.01\% & 99.47\% & 2.52\% \\
  7  & 71.91\% & 99.20\% & 0.05\% \\
  8  & 69.87\% & 99.81\% & 32.60\% \\
  9  & 68.72\% & 98.73\% & 11.04\% \\
  10 & 72.71\% & 99.60\% & 1.02\% \\
  11 & 72.47\% & 99.91\% & 1.03\% \\
  12 & 79.33\% & 99.94\% & 1.25\% \\
  13 & 76.55\% & 99.79\% & 2.36\% \\
  14 & 69.43\% & 99.91\% & 8.33\% \\
  15 & 69.26\% & 99.95\% & 41.25\% \\

  \bottomrule
  \end{tblr}
\end{table}

% \xsubsection{引用公式}{Ref equations}

% 引用公式使用 \clist{\eqref{label}} 命令，将引用公式的编号（带括号）：

% \begin{texcode}[]{}
%   如式 \eqref{eqn:c1:diff} 所示
% \end{texcode}

% \xsubsection{引用其他内容}{Ref other labels}

基于 \clist{\label{}} 编号的引用，可以通过 \clist[columbiablue]{cleveref} 提供的 \clist{\cref} 实现，它可以通过识别 label 的类型，自动为引用添加前缀。具体的使用方法可以查看该宏包的文档。

应当这样使用 \clist{\label{TYPE:NAME}}，其中 TYPE 指此标签的类型，应当使用模板提供的或自行指定；NAME 为此标签的名称，可以自行定义，不应当重复。模板为所有举例过的环境预设了类型，其中在\cref{section:Environments for math}中预设了 $12$种数学环境，则其 TYPE 与该环境的「环境」列保持一致，其余类型的定义见下：

\begin{texcode}[]{}
  如\cref{equation:c1:diff} 所示 \\
  如\crefrange{equation:c1:exp}{equation:c1:diff} 所示 \\
  如\cref{equation:c1:diff,equation:c1:exp,equation:c1:mdl:cstr_dis} 所示 \\
  如\cref{figure:icon} 所示，如\cref{subfigure:icon} 所示 \\
  如\cref{table:methodcompare} 所示 \\
  如\cref{algorithm:SOH} 所示 \\
  如\cref{listing:a_label} 所示 \\
  如\cref{theorem:pt} 所示
\end{texcode}

模板只对数学公式设置了多种引用时的样式，并未对其他类型设置，如果有使用的需要，请自行查看手册与模板文件模仿修改。

\xsection{参考文献}{References}

\begin{tcolorbox}[colback=red!5!white,colframe=red!75!black]
  \begin{enumerate}[leftmargin=0.5cm]
    \item 参考文献的著录格式应符合国家标准GB/T 7714-2015《文后参考文献著录规则》。参考文献中每条项目应齐全。
    \item 参考文献里面标点符号：英文文献用半角,中文文献用全角。
  \end{enumerate}
\end{tcolorbox}

\begin{tcolorbox}[colback=blue!5!white,colframe=blue!75!black]
  \begin{enumerate}[leftmargin=0.5cm]
    \item 文后著录的参考文献务必实事求是。论文中引用过的文献必须著录，未引用的文献不得出现。应遵循学术道德规范，避免涉嫌抄袭、剽窃等学术不端行为。
    \item 参考文献一般应是作者亲自考察过的对学位论文有参考价值的文献，除特殊情况外，一般不应间接引用。
    \item 参考文献应有权威性，要注意引用最新的文献。
    \item 参考文献的数量：
    博士学位论文，一般应在80篇以上，其中，期刊文献60篇以上，国外文献30篇以上，均以近5年的文献为主。
    硕士学位论文，一般应在30篇以上，其中，期刊文献不少于20篇，国外文献不少于10篇，均以近5年的文献为主。
  \end{enumerate}
\end{tcolorbox}

模板使用 biblatex 编译参考文献，默认采用顺序编码制，有多种引用参考文献的方式，请自行查阅 biblatex-gbt-7715-2015 宏包，即 \clist{texdoc biblatex-gbt-7714-2015}。

在导言区使用\clist{\addbibresource{path/to/.bib file}}添加参考文献数据库，默认参考文献数据库位于 ./References/ 下。

这里主要举例三种引用方式，使用 \clist{\cite{ref}} 的引用结果作为上标出现在正文中，使用 \clist{\parencite{ref}} 的引用结果作为正文内容（一个名词）出现在正文中，使用 \clist{\footfullcite{ref}} 的引用结果作为脚注内容在脚注部分出现。



\begin{texcode}[]{}
  这是一个比较常见的问题 \cite{barella_situ_2021}。可以发现 xxxx \cite{atta_enhanced_2021}。\\
  文献\parencite{张燕2013电气自动化在电气工程中的应用探讨} 提出了 xxx，文献 \parencite{黄雪芳2012探讨电气工程中自动化技术的应用} 则提出了 yyy。
\end{texcode}

\footnote{这是脚注1}根据 zzz 发布的报告 \footfullcite{王磊2011电气自动化控制设备可靠性探究}，可以发现\footnote{这是脚注2}